\documentclass[instructions.tex]{subfiles}
\begin{document}
 
\chapter{Jésus souffrant nous apprend à souffrir.}
 
Deux sortes de personnes sont dans l’erreur au sujet des afflictions. Je n’ai point de mérite à souffrir, disent les uns ; j’y sens trop de répugnance. Je souffrirais volontiers, disent les autres, si je l’avais mérité ; mais je suis innocent. Instruisons-les par l’exemple de Jésus souffrant.
 
\primo Quant aux premiers, le Sauveur leur apprend que la répugnance, la sensibilité, les sentimens de la nature, ne sont point contraires à la patience, et n’ôtent pas le mérite, pourvu que dans le fond de l’âme on soit soumis à Dieu. Les consolations qui accompagnent certaines croix n’en font pas le mérite ; elles sont quelquefois accordées à des âmes foibles, comme le lait et le miel qu’on donne aux enfans. Mais les croix qui sont sans consolations et qui accablent, sont comme la nourriture des âmes fortes.
 
Qui fut jamais plus accablé d’ennui, qui eut jamais plus de répugnance à souffrir, que le Sauveur au jardin des Oliviers, lorsque, succombant sous le poids d’une tristesse mortelle, il disait à Dieu son Père : Ah ! mon Père, éloignez de moi ce calice de douleurs ! Quelle fut même sa désolation sur la croix, lorsqu’il s’écriait : Mon Père, m’avez-vous donc abandonné ?
 
Le Sauveur ressentait vivement ses douleurs, parce qu’il était homme ; il les ressentait encore plus vivement, parce qu’il était innocent. Il s’offrit cependant à la croix ; et malgré la répugnance de la nature, malgré la confusion dont il était accablé, son âme sainte fut toujours si parfaitement soumise à la volonté de son Père, qu’il souffrit avec joie la rigueur et l’opprobre de ce supplice\footnote{(1) Proposito sibi gaudio, sustinuit crucem, confusione contempta. Hebr. 12.}.
 
N’a-t-on pas vu des saints à qui la douleur faisait jeter des larmes et pousser les hauts cris, qui néanmoins, dans le fond de l’ame, étaient soumis aux ordres du ciel? Ce n’est donc pas dans l’insensibilité que consiste la patience, mais dans un attachement sincère à la volonté de Dieu.
 
\secundo Quant à vous qui dites qu’il est dur de souffrir la persécution et la calomnie, lorsqu’on est innocent, jetez les yeux sur votre Sauveur, qui est plus innocent que vous : il est votre Dieu, le Saint des saints ; néanmoins il souffre, et jamais il n’y a eu de douleur semblable à la sienne. De quoi donc vous plaindriez-vous? Si votre Dieu souffre, pourquoi ne souffririez-vous pas? Il est innocent, et vous êtes coupable : et fussiez-vous innocent aux yeux des hommes en une chose, l’êtes-vous en tout devant Dieu?
 
Saint Pierre, martyr, accusé faussement, et condamné, quoique innocent, à une dure prison, se plaignait amoureusement au pied de son crucifix : Eh! mon Sauveur, qu’ai-je donc fait pour être traité de la sorte? Le Sauveur lui répondit : Et moi, Pierre, qu’ai-je fait pour être attaché à la croix? Cette parole lui inspira tant de courage, que les peines de tout l’univers n’eussent pas été capables de l’ébranler.
 
On ne vous a point encore noirci par des calomnies aussi atroces que celles dont on a chargé Jésus-Christ. On ne vous a pas garrotté et flagellé comme lui. Vous n’avez pas été attaché à un gibet comme ce Sauveur innocent. Vous n’avez même jamais ressenti ces grandes afflictions dont Dieu visite ses saints : et cependant, pour une légère perte, pour une parole fâcheuse, pour une contradiction, vous êtes déconcerté. Oh! que vous ne comprenez guère de quel esprit vous devez être animé !
 
Plus vous êtes innocent, plus vous êtes heureux et glorieux de souffrir, dit saint Pierre. Quelle gloire auriez-vous de souffrir l’opprobre, lorsque vous l’auriez mérité par un crime\footnote{(1) Quæ enim est gloria, si peccantes et colaphizati suffertis? 1. Petr. 2.}? Plus vous êtes innocent, plus vous aurez de consolations à souffrir et à pardonner, parce que vous serez plus semblable à Jésus-Christ. Il souffre la mort, et il offre le pardon à ses bourreaux ; il embrasse même Judas, qui le trahit ; et vous ne souffririez pas une injure d’un ennemi ou d’un frère ! Apprenez, par l’exemple d’un Dieu, qu’il est bien plus avantageux de recevoir une injure que de la faire ; plus honorable de souffrir que de faire souffrir les autres, et qu’il vaut mieux souffrir de tous les hommes, que d’en faire souffrir ou que d’en haïr un seul.
 
Armez-vous de cette pensée, dit saint Pierre, que Jésus-Christ a souffert pour vous. En jetant les yeux sur ce divin original, vous trouverez dans son innocence un courage à l’épreuve de toutes les contradictions de la vie. Si le Sauveur, montant au Calvaire, vous avait prié de l’aider à porter sa croix, avec quelle ardeur ne l’auriez-vous pas embrassée ! Quoi, vous n’auriez pas balancé de porter cette pesante croix sous laquelle un Homme-Dieu succombe, et vous ne pouvez vous résoudre à porter les plus légères ! Comment souffririez-vous les insultes des bourreaux, vous qui ne pouvez souffrir une parole? Comment vous laisseriez-vous charger de plaies, vous qui ne pouvez souffrir un mépris, une contradiction? Vous auriez eu le courage de porter la croix du Sauveur ; pourquoi donc n’acceptez-vous pas une portion de cette croix dont il vous honore tous les jours? Le serviteur doit-il être plus délicat que son seigneur, et tenir une autre route que celle qu’il voit tenir à son maître\footnote{(1) Non est discipulus super magistrum. Matth. 10. 24.}?
 
Finissons par cette pensée de saint Bonaventure, qui enseigne que Jésus-Christ n’a pas été un seul moment sans souffrir ; que, dès le premier instant de sa conception jusqu’à sa mort, il ressentit en son âme les horreurs de la croix ; qu’il la désira, et que, les tourmens de sa passion furent toujours présens à mon esprit. Ce saint docteur assure même que le Fils de Dieu était disposé à s’incarner, à être attaché à la croix dès le commencement du monde, et à y souffrir les mépris, son argent, sont comme une boue épaisse qu’il amasse autour de lui, dans laquelle il ensevelit son cœur\footnote{(1) Melins est modicum justo, super divitias peccatorum multas. Ps. 36. -- (2) Radix enim omnium malorum est cupiditas. 1. Tim. 6.}.
 
Oh! qu’on est insensé d’attacher ainsi son cœur à la terre, et d’être esclave de sa propre cupidité !
 
A l’exemple de Salomon, ne demandez ni les richesses, ni la pauvreté. Si vous êtes pauvre, ne désirez pas d’être riche. Ces désirs, dit saint Paul, sont inutiles, nuisibles, et plongent l’homme dans la mort. Si vous êtes riche, modérez vos désirs : on est toujours malheureux quand on n’est pas content de ce qu’on possède.
 
\section{Résolutions}
 
\primo Persuadé que les richesses ne rendent pas heureux celui qui les possède, qu’on n’en jouit pas long-temps, puisqu’on ne les emporte pas en l’autre vie, et qu’elles peuvent devenir aisément une source féconde de péchés, et la cause de la damnation éternelle, je ne m’y attacherai point.

\secundo J’en supporterai avec résignation la perte.

\tertio Et je pardonnerai volontiers les torts qui me seront faits.

\prayer{O mon Dieu! ne permettez pas que mon cœur que vous avez formé pour vous, s’attache à d’autres objets au préjudice de l’amour que je vous dois.}
 
\end{document}
